\section{RISC-V TEE lineage: Sanctum / Keystone / Penglai / SPEAR-V}

\begin{frame}[t,shrink=6]{RISC-V TEE lineage: Sanctum / Keystone / Penglai / SPEAR-V：方向开场}
\DeckKicker{方向开场}
\SlideMeta{证据等级}{方向综述}{来源状态：公开材料综合}
{\large\bfseries\color{Ink} RISC-V TEE lineage: Sanctum / Keystone / Penglai / SPEAR-V 的核心是先界定保护对象、管理边界和证据强度，再用三篇主讲材料串起技术演进。\par}\vspace{1.8mm}
\begin{columns}[T,totalwidth=\textwidth]
  \begin{column}{0.45\textwidth}
    \fcolorbox{Rule}{Paper}{%
  \begin{minipage}[t]{0.97\linewidth}
    \vspace{1.1mm}
    \hspace{1.4mm}\begin{minipage}{0.92\linewidth}
      \PanelTitle{内容说明}
      \scriptsize \begin{itemize}
  \item 从 PMP-based enclave 到 scalable memory protection，再到 tag/metadata primitive。
  \item 固定三篇主讲材料；不引入辅助材料，直接呈现 RISC-V enclave lineage。
  \item 先讲方向痛点和设计轴，再逐篇说明贡献、限制、证据强度和可落地条件。
  \item 对规范、Survey、draft、vendor evidence 保持显式标注，避免把背景材料写成一手系统证明。
\end{itemize}
    \end{minipage}
    \vspace{1mm}
  \end{minipage}}
  \end{column}
  \begin{column}{0.49\textwidth}
    \fcolorbox{Rule}{Panel}{%
  \begin{minipage}[t]{0.97\linewidth}
    \vspace{1.1mm}
    \hspace{1.4mm}\begin{minipage}{0.92\linewidth}
      \PanelTitle{三篇主讲材料的分工}
      \scriptsize {\tiny
\begin{tabularx}{\linewidth}{@{}YYY@{}}
\toprule
\textbf{论文} & \textbf{定位} & \textbf{证据边界} \\
\midrule
Keystone: 开源 RISC-V TEE 框架 & 开创/基础入口 & E1 / 基础证据 \\
Penglai: scalable enclave memory protection & 代表性改进 & E1 / 同行评审改进证据 \\
SPEAR-V: low-overhead RISC-V enclave primitive & 当前 SOTA/补充边界 & E1 / 同行评审改进证据 \\
\bottomrule
\end{tabularx}
}
    \end{minipage}
    \vspace{1mm}
  \end{minipage}}
  \end{column}
\end{columns}
\vspace{1mm}
{\tiny\textcolor{Muted}{引用依据：本页综合三篇主讲材料的研究定位、证据等级和公开来源状态。}}
\end{frame}


\begin{frame}[t,shrink=6]{Keystone: 开源 RISC-V TEE 框架：内容摘要}
\DeckKicker{内容摘要}
\SlideMeta{证据等级}{E1 / 基础证据}{来源状态：作者公开 PDF 已核验}
{\large\bfseries\color{Ink} 开源 RISC-V TEE 框架，用 PMP、security monitor 和 runtime abstraction 构建 enclave\par}\vspace{1.8mm}
\begin{columns}[T,totalwidth=\textwidth]
  \begin{column}{0.45\textwidth}
    \fcolorbox{Rule}{Paper}{%
  \begin{minipage}[t]{0.97\linewidth}
    \vspace{1.1mm}
    \hspace{1.4mm}\begin{minipage}{0.92\linewidth}
      \PanelTitle{内容说明}
      \scriptsize \begin{itemize}
  \item 开源 RISC-V TEE 框架，用 PMP、security monitor 和 runtime abstraction 构建 enclave
  \item EuroSys 2020 Keystone is the foundational public RISC-V open-source TEE framework paper.
\end{itemize}
    \end{minipage}
    \vspace{1mm}
  \end{minipage}}
  \end{column}
  \begin{column}{0.49\textwidth}
    \fcolorbox{Rule}{Panel}{%
  \begin{minipage}[t]{0.97\linewidth}
    \vspace{1.1mm}
    \hspace{1.4mm}\begin{minipage}{0.92\linewidth}
      \PanelTitle{证据定位}
      \scriptsize {\tiny
\begin{tabularx}{\linewidth}{@{}YY@{}}
\toprule
\textbf{字段} & \textbf{内容} \\
\midrule
论文定位 & 开创/基础入口 \\
材料类型 & 系统论文 \\
证据等级 & E1 / 基础证据 \\
来源状态 & 作者公开 PDF 已核验 \\
\bottomrule
\end{tabularx}
}
    \end{minipage}
    \vspace{1mm}
  \end{minipage}}
  \end{column}
\end{columns}
\vspace{1mm}
{\tiny\textcolor{Muted}{引用依据：引用依据为论文原文、官方规范或公开材料；本页结论按 E1 / 基础证据 的证据等级使用，不外推到未验证平台。}}
\end{frame}


\begin{frame}[t,shrink=6]{Keystone: 开源 RISC-V TEE 框架：研究背景}
\DeckKicker{研究背景}
\SlideMeta{证据等级}{E1 / 基础证据}{来源状态：作者公开 PDF 已核验}
{\large\bfseries\color{Ink} 商业 TEE 难定制、难验证\par}\vspace{1.8mm}
\begin{columns}[T,totalwidth=\textwidth]
  \begin{column}{0.45\textwidth}
    \fcolorbox{Rule}{Paper}{%
  \begin{minipage}[t]{0.97\linewidth}
    \vspace{1.1mm}
    \hspace{1.4mm}\begin{minipage}{0.92\linewidth}
      \PanelTitle{内容说明}
      \scriptsize \begin{itemize}
  \item 商业 TEE 难定制、难验证
  \item RISC-V 提供开放硬件/软件协同机会，需要可研究、可裁剪的 enclave baseline
\end{itemize}
    \end{minipage}
    \vspace{1mm}
  \end{minipage}}
  \end{column}
  \begin{column}{0.49\textwidth}
    \fcolorbox{Rule}{Panel}{%
  \begin{minipage}[t]{0.97\linewidth}
    \vspace{1.1mm}
    \hspace{1.4mm}\begin{minipage}{0.92\linewidth}
      \PanelTitle{问题边界}
      \scriptsize {\tiny
\begin{tabularx}{\linewidth}{@{}YY@{}}
\toprule
\textbf{维度} & \textbf{说明} \\
\midrule
保护对象 & RISC-V TEE lineage: Sanctum / Keystone / Penglai / SPEAR-V \\
传统瓶颈 & 从 PMP-based enclave 到 scalable memory protection，再到 tag/metadata prim。 \\
本文入口 & Keystone: 开源 RISC-V TEE 框架 \\
证据限制 & E1 / 基础证据 \\
\bottomrule
\end{tabularx}
}
    \end{minipage}
    \vspace{1mm}
  \end{minipage}}
  \end{column}
\end{columns}
\vspace{1mm}
{\tiny\textcolor{Muted}{引用依据：引用依据为论文原文、官方规范或公开材料；本页结论按 E1 / 基础证据 的证据等级使用，不外推到未验证平台。}}
\end{frame}


\begin{frame}[t,shrink=6]{Keystone: 开源 RISC-V TEE 框架：解决方案}
\DeckKicker{解决方案}
\SlideMeta{证据等级}{E1 / 基础证据}{来源状态：作者公开 PDF 已核验}
{\large\bfseries\color{Ink} 把 PMP region、security monitor、runtime 和 host ABI 组合成可扩展的开源 TEE 框架\par}\vspace{1.8mm}
\begin{columns}[T,totalwidth=\textwidth]
  \begin{column}{0.45\textwidth}
    \fcolorbox{Rule}{Paper}{%
  \begin{minipage}[t]{0.97\linewidth}
    \vspace{1.1mm}
    \hspace{1.4mm}\begin{minipage}{0.92\linewidth}
      \PanelTitle{内容说明}
      \scriptsize \begin{itemize}
  \item 把 PMP region、security monitor、runtime 和 host ABI 组合成可扩展的开源 TEE 框架
\end{itemize}
    \end{minipage}
    \vspace{1mm}
  \end{minipage}}
  \end{column}
  \begin{column}{0.49\textwidth}
    \fcolorbox{Rule}{Panel}{%
  \begin{minipage}[t]{0.97\linewidth}
    \vspace{1.1mm}
    \hspace{1.4mm}\begin{minipage}{0.92\linewidth}
      \PanelTitle{核心设计拆解}
      \scriptsize \begin{itemize}
  \item 01. PMP-based enclave isolation
  \item 02. Security monitor and runtime abstraction
  \item 03. Open RISC-V TEE software/hardware co-design
\end{itemize}
    \end{minipage}
    \vspace{1mm}
  \end{minipage}}
  \end{column}
\end{columns}
\vspace{1mm}
{\tiny\textcolor{Muted}{引用依据：引用依据为论文原文、官方规范或公开材料；本页结论按 E1 / 基础证据 的证据等级使用，不外推到未验证平台。}}
\end{frame}


\begin{frame}[t,shrink=6]{Keystone: 开源 RISC-V TEE 框架：实验结果}
\DeckKicker{实验结果}
\SlideMeta{证据等级}{E1 / 基础证据}{来源状态：作者公开 PDF 已核验}
{\large\bfseries\color{Ink} 论文报告 CoreMark/Beebs/RV8 等开销结果\par}\vspace{1.8mm}
\begin{columns}[T,totalwidth=\textwidth]
  \begin{column}{0.45\textwidth}
    \fcolorbox{Rule}{Paper}{%
  \begin{minipage}[t]{0.97\linewidth}
    \vspace{1.1mm}
    \hspace{1.4mm}\begin{minipage}{0.92\linewidth}
      \PanelTitle{内容说明}
      \scriptsize \begin{itemize}
  \item 论文报告 CoreMark/Beebs/RV8 等开销结果
  \item 具体数字以本地 PDF 为准
\end{itemize}
    \end{minipage}
    \vspace{1mm}
  \end{minipage}}
  \end{column}
  \begin{column}{0.49\textwidth}
    \fcolorbox{Rule}{Panel}{%
  \begin{minipage}[t]{0.97\linewidth}
    \vspace{1.1mm}
    \hspace{1.4mm}\begin{minipage}{0.92\linewidth}
      \PanelTitle{实验/证据边界}
      \scriptsize {\tiny
\begin{tabularx}{\linewidth}{@{}YYY@{}}
\toprule
\textbf{对象} & \textbf{可支撑} & \textbf{边界} \\
\midrule
数据/来源 & Keystone: 开源 RISC-V TEE 框架 & 作者公开 PDF 已核验 \\
Baseline/对照 & 以论文或规范原文为准 & 不跨材料外推 \\
指标/对象 & 只使用当前来源可证明的信息 & E1 / 基础证据 \\
使用边界 & 可进入本方向演进图 & 不能替代更强证据 \\
\bottomrule
\end{tabularx}
}
    \end{minipage}
    \vspace{1mm}
  \end{minipage}}
  \end{column}
\end{columns}
\vspace{1mm}
{\tiny\textcolor{Muted}{引用依据：引用依据为论文原文、官方规范或公开材料；本页结论按 E1 / 基础证据 的证据等级使用，不外推到未验证平台。}}
\end{frame}


\begin{frame}[t,shrink=6]{Keystone: 开源 RISC-V TEE 框架：文章评价}
\DeckKicker{文章评价}
\SlideMeta{证据等级}{E1 / 基础证据}{来源状态：作者公开 PDF 已核验}
{\large\bfseries\color{Ink} RISC-V TEE baseline，适合解释 PMP-based enclave 的优势和限制\par}\vspace{1.8mm}
\begin{columns}[T,totalwidth=\textwidth]
  \begin{column}{0.45\textwidth}
    \fcolorbox{Rule}{Paper}{%
  \begin{minipage}[t]{0.97\linewidth}
    \vspace{1.1mm}
    \hspace{1.4mm}\begin{minipage}{0.92\linewidth}
      \PanelTitle{内容说明}
      \scriptsize \begin{itemize}
  \item 设计优势：RISC-V TEE baseline，适合解释 PMP-based enclave 的优势和限制。
  \item 局限性：EuroSys 2020 foundational evidence，不代表后续标准化 confidential VM 或生产 enclave。
  \item 商业化潜力：开源可定制性强；风险在 PMP region 限制、生态成熟度和安全 monitor 正确性。
  \item 证据边界：E1 / Foundational；作者公开 PDF 已核验。
\end{itemize}
    \end{minipage}
    \vspace{1mm}
  \end{minipage}}
  \end{column}
  \begin{column}{0.49\textwidth}
    \fcolorbox{Rule}{Panel}{%
  \begin{minipage}[t]{0.97\linewidth}
    \vspace{1.1mm}
    \hspace{1.4mm}\begin{minipage}{0.92\linewidth}
      \PanelTitle{文章评价}
      \scriptsize {\tiny
\begin{tabularx}{\linewidth}{@{}YY@{}}
\toprule
\textbf{维度} & \textbf{判断} \\
\midrule
设计优势 & RISC-V TEE baseline，适合解释 PMP-based enclave 的优势和限制。 \\
主要局限 & EuroSys 2020 foundational evidence，不代表后续标准化 confidential VM 或生产 enclave。 \\
商业落地 & 开源可定制性强；风险在 PMP region 限制、生态成熟度和安全 monitor 正确性。 \\
讲稿定位 & 开创/基础入口; 证据强度 E1 \\
\bottomrule
\end{tabularx}
}
    \end{minipage}
    \vspace{1mm}
  \end{minipage}}
  \end{column}
\end{columns}
\vspace{1mm}
{\tiny\textcolor{Muted}{引用依据：引用依据为论文原文、官方规范或公开材料；本页结论按 E1 / 基础证据 的证据等级使用，不外推到未验证平台。}}
\end{frame}


\begin{frame}[t,shrink=6]{Penglai: scalable enclave memory protection：内容摘要}
\DeckKicker{内容摘要}
\SlideMeta{证据等级}{E1 / 同行评审改进证据}{来源状态：本地 PDF 已核验}
{\large\bfseries\color{Ink} 面向大规模 enclave 的 scalable memory protection，回应 serverless/microservice 对大量。\par}\vspace{1.8mm}
\begin{columns}[T,totalwidth=\textwidth]
  \begin{column}{0.45\textwidth}
    \fcolorbox{Rule}{Paper}{%
  \begin{minipage}[t]{0.97\linewidth}
    \vspace{1.1mm}
    \hspace{1.4mm}\begin{minipage}{0.92\linewidth}
      \PanelTitle{内容说明}
      \scriptsize \begin{itemize}
  \item 面向大规模 enclave 的 scalable memory protection，回应 serverless/microservice 对大量动态 enclave 的需求
  \item OSDI 2021 Penglai is a peer-reviewed SOTA scalable RISC-V enclave OS after Keystone.
\end{itemize}
    \end{minipage}
    \vspace{1mm}
  \end{minipage}}
  \end{column}
  \begin{column}{0.49\textwidth}
    \fcolorbox{Rule}{Panel}{%
  \begin{minipage}[t]{0.97\linewidth}
    \vspace{1.1mm}
    \hspace{1.4mm}\begin{minipage}{0.92\linewidth}
      \PanelTitle{证据定位}
      \scriptsize {\tiny
\begin{tabularx}{\linewidth}{@{}YY@{}}
\toprule
\textbf{字段} & \textbf{内容} \\
\midrule
论文定位 & 代表性改进 \\
材料类型 & 系统论文 \\
证据等级 & E1 / 同行评审改进证据 \\
来源状态 & 本地 PDF 已核验 \\
\bottomrule
\end{tabularx}
}
    \end{minipage}
    \vspace{1mm}
  \end{minipage}}
  \end{column}
\end{columns}
\vspace{1mm}
{\tiny\textcolor{Muted}{引用依据：引用依据为论文原文、官方规范或公开材料；本页结论按 E1 / 同行评审改进证据 的证据等级使用，不外推到未验证平台。}}
\end{frame}


\begin{frame}[t,shrink=6]{Penglai: scalable enclave memory protection：研究背景}
\DeckKicker{研究背景}
\SlideMeta{证据等级}{E1 / 同行评审改进证据}{来源状态：本地 PDF 已核验}
{\large\bfseries\color{Ink} 传统 PMP region 数量有限，难以支撑大量 enclave、动态生命周期和更大的受保护内存\par}\vspace{1.8mm}
\begin{columns}[T,totalwidth=\textwidth]
  \begin{column}{0.45\textwidth}
    \fcolorbox{Rule}{Paper}{%
  \begin{minipage}[t]{0.97\linewidth}
    \vspace{1.1mm}
    \hspace{1.4mm}\begin{minipage}{0.92\linewidth}
      \PanelTitle{内容说明}
      \scriptsize \begin{itemize}
  \item 传统 PMP region 数量有限，难以支撑大量 enclave、动态生命周期和更大的受保护内存
\end{itemize}
    \end{minipage}
    \vspace{1mm}
  \end{minipage}}
  \end{column}
  \begin{column}{0.49\textwidth}
    \fcolorbox{Rule}{Panel}{%
  \begin{minipage}[t]{0.97\linewidth}
    \vspace{1.1mm}
    \hspace{1.4mm}\begin{minipage}{0.92\linewidth}
      \PanelTitle{问题边界}
      \scriptsize {\tiny
\begin{tabularx}{\linewidth}{@{}YY@{}}
\toprule
\textbf{维度} & \textbf{说明} \\
\midrule
保护对象 & RISC-V TEE lineage: Sanctum / Keystone / Penglai / SPEAR-V \\
传统瓶颈 & 从 PMP-based enclave 到 scalable memory protection，再到 tag/metadata prim。 \\
本文入口 & Penglai: scalable enclave memory protection \\
证据限制 & E1 / 同行评审改进证据 \\
\bottomrule
\end{tabularx}
}
    \end{minipage}
    \vspace{1mm}
  \end{minipage}}
  \end{column}
\end{columns}
\vspace{1mm}
{\tiny\textcolor{Muted}{引用依据：引用依据为论文原文、官方规范或公开材料；本页结论按 E1 / 同行评审改进证据 的证据等级使用，不外推到未验证平台。}}
\end{frame}


\begin{frame}[t,shrink=6]{Penglai: scalable enclave memory protection：解决方案}
\DeckKicker{解决方案}
\SlideMeta{证据等级}{E1 / 同行评审改进证据}{来源状态：本地 PDF 已核验}
{\large\bfseries\color{Ink} 引入更可扩展的内存保护和 metadata 管理机制，把 RISC-V enclave 从少量静态隔离推进到云端规模化\par}\vspace{1.8mm}
\begin{columns}[T,totalwidth=\textwidth]
  \begin{column}{0.45\textwidth}
    \fcolorbox{Rule}{Paper}{%
  \begin{minipage}[t]{0.97\linewidth}
    \vspace{1.1mm}
    \hspace{1.4mm}\begin{minipage}{0.92\linewidth}
      \PanelTitle{内容说明}
      \scriptsize \begin{itemize}
  \item 引入更可扩展的内存保护和 metadata 管理机制，把 RISC-V enclave 从少量静态隔离推进到云端规模化
\end{itemize}
    \end{minipage}
    \vspace{1mm}
  \end{minipage}}
  \end{column}
  \begin{column}{0.49\textwidth}
    \fcolorbox{Rule}{Panel}{%
  \begin{minipage}[t]{0.97\linewidth}
    \vspace{1.1mm}
    \hspace{1.4mm}\begin{minipage}{0.92\linewidth}
      \PanelTitle{核心设计拆解}
      \scriptsize \begin{itemize}
  \item 01. Scalable enclave memory protection
  \item 02. Metadata management for enclave lifecycle
  \item 03. Serverless/microservice-oriented enclave scaling
\end{itemize}
    \end{minipage}
    \vspace{1mm}
  \end{minipage}}
  \end{column}
\end{columns}
\vspace{1mm}
{\tiny\textcolor{Muted}{引用依据：引用依据为论文原文、官方规范或公开材料；本页结论按 E1 / 同行评审改进证据 的证据等级使用，不外推到未验证平台。}}
\end{frame}


\begin{frame}[t,shrink=6]{Penglai: scalable enclave memory protection：实验结果}
\DeckKicker{实验结果}
\SlideMeta{证据等级}{E1 / 同行评审改进证据}{来源状态：本地 PDF 已核验}
{\large\bfseries\color{Ink} OSDI 论文报告可支持大量 enclave 和较大安全内存\par}\vspace{1.8mm}
\begin{columns}[T,totalwidth=\textwidth]
  \begin{column}{0.45\textwidth}
    \fcolorbox{Rule}{Paper}{%
  \begin{minipage}[t]{0.97\linewidth}
    \vspace{1.1mm}
    \hspace{1.4mm}\begin{minipage}{0.92\linewidth}
      \PanelTitle{内容说明}
      \scriptsize \begin{itemize}
  \item OSDI 论文报告可支持大量 enclave 和较大安全内存
  \item 具体指标以本地 PDF 为准
\end{itemize}
    \end{minipage}
    \vspace{1mm}
  \end{minipage}}
  \end{column}
  \begin{column}{0.49\textwidth}
    \fcolorbox{Rule}{Panel}{%
  \begin{minipage}[t]{0.97\linewidth}
    \vspace{1.1mm}
    \hspace{1.4mm}\begin{minipage}{0.92\linewidth}
      \PanelTitle{实验/证据边界}
      \scriptsize {\tiny
\begin{tabularx}{\linewidth}{@{}YYY@{}}
\toprule
\textbf{对象} & \textbf{可支撑} & \textbf{边界} \\
\midrule
数据/来源 & Penglai: scalable enclave memory protection & 本地 PDF 已核验 \\
Baseline/对照 & 以论文或规范原文为准 & 不跨材料外推 \\
指标/对象 & 只使用当前来源可证明的信息 & E1 / 同行评审改进证据 \\
使用边界 & 可进入本方向演进图 & 不能替代更强证据 \\
\bottomrule
\end{tabularx}
}
    \end{minipage}
    \vspace{1mm}
  \end{minipage}}
  \end{column}
\end{columns}
\vspace{1mm}
{\tiny\textcolor{Muted}{引用依据：引用依据为论文原文、官方规范或公开材料；本页结论按 E1 / 同行评审改进证据 的证据等级使用，不外推到未验证平台。}}
\end{frame}


\begin{frame}[t,shrink=6]{Penglai: scalable enclave memory protection：文章评价}
\DeckKicker{文章评价}
\SlideMeta{证据等级}{E1 / 同行评审改进证据}{来源状态：本地 PDF 已核验}
{\large\bfseries\color{Ink} 系统贡献强，直接解决 PMP-based enclave 在规模上的瓶颈\par}\vspace{1.8mm}
\begin{columns}[T,totalwidth=\textwidth]
  \begin{column}{0.45\textwidth}
    \fcolorbox{Rule}{Paper}{%
  \begin{minipage}[t]{0.97\linewidth}
    \vspace{1.1mm}
    \hspace{1.4mm}\begin{minipage}{0.92\linewidth}
      \PanelTitle{内容说明}
      \scriptsize \begin{itemize}
  \item 设计优势：系统贡献强，直接解决 PMP-based enclave 在规模上的瓶颈。
  \item 局限性：需要硬件修改和生态配套，不等同于 RISC-V ratified confidential-computing standard。
  \item 商业化潜力：适合云原生 enclave 服务方向；风险在硬件采纳、标准化和云平台管理接口。
  \item 证据边界：E1 / Peer-reviewed SOTA；本地 PDF 已核验。
\end{itemize}
    \end{minipage}
    \vspace{1mm}
  \end{minipage}}
  \end{column}
  \begin{column}{0.49\textwidth}
    \fcolorbox{Rule}{Panel}{%
  \begin{minipage}[t]{0.97\linewidth}
    \vspace{1.1mm}
    \hspace{1.4mm}\begin{minipage}{0.92\linewidth}
      \PanelTitle{文章评价}
      \scriptsize {\tiny
\begin{tabularx}{\linewidth}{@{}YY@{}}
\toprule
\textbf{维度} & \textbf{判断} \\
\midrule
设计优势 & 系统贡献强，直接解决 PMP-based enclave 在规模上的瓶颈。 \\
主要局限 & 需要硬件修改和生态配套，不等同于 RISC-V ratified confidential-computing standard。 \\
商业落地 & 适合云原生 enclave 服务方向；风险在硬件采纳、标准化和云平台管理接口。 \\
讲稿定位 & 代表性改进; 证据强度 E1 \\
\bottomrule
\end{tabularx}
}
    \end{minipage}
    \vspace{1mm}
  \end{minipage}}
  \end{column}
\end{columns}
\vspace{1mm}
{\tiny\textcolor{Muted}{引用依据：引用依据为论文原文、官方规范或公开材料；本页结论按 E1 / 同行评审改进证据 的证据等级使用，不外推到未验证平台。}}
\end{frame}


\begin{frame}[t,shrink=6]{SPEAR-V: low-overhead RISC-V enclave primitive：内容摘要}
\DeckKicker{内容摘要}
\SlideMeta{证据等级}{E1 / 同行评审改进证据}{来源状态：本地 PDF 已核验}
{\large\bfseries\color{Ink} 提出低开销、实用的 RISC-V enclave primitive，用 tag/metadata primitive 支持双向 sandbox。\par}\vspace{1.8mm}
\begin{columns}[T,totalwidth=\textwidth]
  \begin{column}{0.45\textwidth}
    \fcolorbox{Rule}{Paper}{%
  \begin{minipage}[t]{0.97\linewidth}
    \vspace{1.1mm}
    \hspace{1.4mm}\begin{minipage}{0.92\linewidth}
      \PanelTitle{内容说明}
      \scriptsize \begin{itemize}
  \item 提出低开销、实用的 RISC-V enclave primitive，用 tag/metadata primitive 支持双向 sandbox 和嵌套隔离
  \item SPEAR-V is a peer-reviewed SOTA RISC-V TEE system with SoC and remote-attestation focus.
\end{itemize}
    \end{minipage}
    \vspace{1mm}
  \end{minipage}}
  \end{column}
  \begin{column}{0.49\textwidth}
    \fcolorbox{Rule}{Panel}{%
  \begin{minipage}[t]{0.97\linewidth}
    \vspace{1.1mm}
    \hspace{1.4mm}\begin{minipage}{0.92\linewidth}
      \PanelTitle{证据定位}
      \scriptsize {\tiny
\begin{tabularx}{\linewidth}{@{}YY@{}}
\toprule
\textbf{字段} & \textbf{内容} \\
\midrule
论文定位 & 当前 SOTA/补充边界 \\
材料类型 & 系统论文 \\
证据等级 & E1 / 同行评审改进证据 \\
来源状态 & 本地 PDF 已核验 \\
\bottomrule
\end{tabularx}
}
    \end{minipage}
    \vspace{1mm}
  \end{minipage}}
  \end{column}
\end{columns}
\vspace{1mm}
{\tiny\textcolor{Muted}{引用依据：引用依据为论文原文、官方规范或公开材料；本页结论按 E1 / 同行评审改进证据 的证据等级使用，不外推到未验证平台。}}
\end{frame}


\begin{frame}[t,shrink=6]{SPEAR-V: low-overhead RISC-V enclave primitive：研究背景}
\DeckKicker{研究背景}
\SlideMeta{证据等级}{E1 / 同行评审改进证据}{来源状态：本地 PDF 已核验}
{\large\bfseries\color{Ink} RISC-V enclave 需要更灵活、低开销、可抵抗部分 controlled-channel 风险的硬件 primitive\par}\vspace{1.8mm}
\begin{columns}[T,totalwidth=\textwidth]
  \begin{column}{0.45\textwidth}
    \fcolorbox{Rule}{Paper}{%
  \begin{minipage}[t]{0.97\linewidth}
    \vspace{1.1mm}
    \hspace{1.4mm}\begin{minipage}{0.92\linewidth}
      \PanelTitle{内容说明}
      \scriptsize \begin{itemize}
  \item RISC-V enclave 需要更灵活、低开销、可抵抗部分 controlled-channel 风险的硬件 primitive
\end{itemize}
    \end{minipage}
    \vspace{1mm}
  \end{minipage}}
  \end{column}
  \begin{column}{0.49\textwidth}
    \fcolorbox{Rule}{Panel}{%
  \begin{minipage}[t]{0.97\linewidth}
    \vspace{1.1mm}
    \hspace{1.4mm}\begin{minipage}{0.92\linewidth}
      \PanelTitle{问题边界}
      \scriptsize {\tiny
\begin{tabularx}{\linewidth}{@{}YY@{}}
\toprule
\textbf{维度} & \textbf{说明} \\
\midrule
保护对象 & RISC-V TEE lineage: Sanctum / Keystone / Penglai / SPEAR-V \\
传统瓶颈 & 从 PMP-based enclave 到 scalable memory protection，再到 tag/metadata prim。 \\
本文入口 & SPEAR-V: low-overhead RISC-V enclave primitive \\
证据限制 & E1 / 同行评审改进证据 \\
\bottomrule
\end{tabularx}
}
    \end{minipage}
    \vspace{1mm}
  \end{minipage}}
  \end{column}
\end{columns}
\vspace{1mm}
{\tiny\textcolor{Muted}{引用依据：引用依据为论文原文、官方规范或公开材料；本页结论按 E1 / 同行评审改进证据 的证据等级使用，不外推到未验证平台。}}
\end{frame}


\begin{frame}[t,shrink=6]{SPEAR-V: low-overhead RISC-V enclave primitive：解决方案}
\DeckKicker{解决方案}
\SlideMeta{证据等级}{E1 / 同行评审改进证据}{来源状态：本地 PDF 已核验}
{\large\bfseries\color{Ink} 用 tag/metadata primitive 替代仅靠 PMP region 的隔离表达，支持更细粒度的 protected/unprotec。\par}\vspace{1.8mm}
\begin{columns}[T,totalwidth=\textwidth]
  \begin{column}{0.45\textwidth}
    \fcolorbox{Rule}{Paper}{%
  \begin{minipage}[t]{0.97\linewidth}
    \vspace{1.1mm}
    \hspace{1.4mm}\begin{minipage}{0.92\linewidth}
      \PanelTitle{内容说明}
      \scriptsize \begin{itemize}
  \item 用 tag/metadata primitive 替代仅靠 PMP region 的隔离表达，支持更细粒度的 protected/unprotected 交互
\end{itemize}
    \end{minipage}
    \vspace{1mm}
  \end{minipage}}
  \end{column}
  \begin{column}{0.49\textwidth}
    \fcolorbox{Rule}{Panel}{%
  \begin{minipage}[t]{0.97\linewidth}
    \vspace{1.1mm}
    \hspace{1.4mm}\begin{minipage}{0.92\linewidth}
      \PanelTitle{核心设计拆解}
      \scriptsize \begin{itemize}
  \item 01. Tag/metadata enclave primitive
  \item 02. Bidirectional sandbox and nested isolation
  \item 03. Low-overhead protected/unprotected transitions
\end{itemize}
    \end{minipage}
    \vspace{1mm}
  \end{minipage}}
  \end{column}
\end{columns}
\vspace{1mm}
{\tiny\textcolor{Muted}{引用依据：引用依据为论文原文、官方规范或公开材料；本页结论按 E1 / 同行评审改进证据 的证据等级使用，不外推到未验证平台。}}
\end{frame}


\begin{frame}[t,shrink=6]{SPEAR-V: low-overhead RISC-V enclave primitive：实验结果}
\DeckKicker{实验结果}
\SlideMeta{证据等级}{E1 / 同行评审改进证据}{来源状态：本地 PDF 已核验}
{\large\bfseries\color{Ink} AsiaCCS 论文报告 protected/unprotected 场景低开销\par}\vspace{1.8mm}
\begin{columns}[T,totalwidth=\textwidth]
  \begin{column}{0.45\textwidth}
    \fcolorbox{Rule}{Paper}{%
  \begin{minipage}[t]{0.97\linewidth}
    \vspace{1.1mm}
    \hspace{1.4mm}\begin{minipage}{0.92\linewidth}
      \PanelTitle{内容说明}
      \scriptsize \begin{itemize}
  \item AsiaCCS 论文报告 protected/unprotected 场景低开销
  \item 具体数值以本地 PDF 为准
\end{itemize}
    \end{minipage}
    \vspace{1mm}
  \end{minipage}}
  \end{column}
  \begin{column}{0.49\textwidth}
    \fcolorbox{Rule}{Panel}{%
  \begin{minipage}[t]{0.97\linewidth}
    \vspace{1.1mm}
    \hspace{1.4mm}\begin{minipage}{0.92\linewidth}
      \PanelTitle{实验/证据边界}
      \scriptsize {\tiny
\begin{tabularx}{\linewidth}{@{}YYY@{}}
\toprule
\textbf{对象} & \textbf{可支撑} & \textbf{边界} \\
\midrule
数据/来源 & SPEAR-V: low-overhead RISC-V enclave primitive & 本地 PDF 已核验 \\
Baseline/对照 & 以论文或规范原文为准 & 不跨材料外推 \\
指标/对象 & 只使用当前来源可证明的信息 & E1 / 同行评审改进证据 \\
使用边界 & 可进入本方向演进图 & 不能替代更强证据 \\
\bottomrule
\end{tabularx}
}
    \end{minipage}
    \vspace{1mm}
  \end{minipage}}
  \end{column}
\end{columns}
\vspace{1mm}
{\tiny\textcolor{Muted}{引用依据：引用依据为论文原文、官方规范或公开材料；本页结论按 E1 / 同行评审改进证据 的证据等级使用，不外推到未验证平台。}}
\end{frame}


\begin{frame}[t,shrink=6]{SPEAR-V: low-overhead RISC-V enclave primitive：文章评价}
\DeckKicker{文章评价}
\SlideMeta{证据等级}{E1 / 同行评审改进证据}{来源状态：本地 PDF 已核验}
{\large\bfseries\color{Ink} 机制设计清晰，展示 RISC-V enclave primitive 可沿 tag/metadata 路线演进\par}\vspace{1.8mm}
\begin{columns}[T,totalwidth=\textwidth]
  \begin{column}{0.45\textwidth}
    \fcolorbox{Rule}{Paper}{%
  \begin{minipage}[t]{0.97\linewidth}
    \vspace{1.1mm}
    \hspace{1.4mm}\begin{minipage}{0.92\linewidth}
      \PanelTitle{内容说明}
      \scriptsize \begin{itemize}
  \item 设计优势：机制设计清晰，展示 RISC-V enclave primitive 可沿 tag/metadata 路线演进。
  \item 局限性：Peer-reviewed 系统证据不等同于 RISC-V 标准或量产硬件采纳。
  \item 商业化潜力：若硬件供应商采纳，适合细粒度隔离和嵌入式/边缘 TEE；风险在 ISA 扩展标准化和工具链支持。
  \item 证据边界：E1 / Peer-reviewed SOTA；本地 PDF 已核验。
\end{itemize}
    \end{minipage}
    \vspace{1mm}
  \end{minipage}}
  \end{column}
  \begin{column}{0.49\textwidth}
    \fcolorbox{Rule}{Panel}{%
  \begin{minipage}[t]{0.97\linewidth}
    \vspace{1.1mm}
    \hspace{1.4mm}\begin{minipage}{0.92\linewidth}
      \PanelTitle{文章评价}
      \scriptsize {\tiny
\begin{tabularx}{\linewidth}{@{}YY@{}}
\toprule
\textbf{维度} & \textbf{判断} \\
\midrule
设计优势 & 机制设计清晰，展示 RISC-V enclave primitive 可沿 tag/metadata 路线演进。 \\
主要局限 & Peer-reviewed 系统证据不等同于 RISC-V 标准或量产硬件采纳。 \\
商业落地 & 若硬件供应商采纳，适合细粒度隔离和嵌入式/边缘 TEE；风险在 ISA 扩展标准化和工具链支持。 \\
讲稿定位 & 当前 SOTA/补充边界; 证据强度 E1 \\
\bottomrule
\end{tabularx}
}
    \end{minipage}
    \vspace{1mm}
  \end{minipage}}
  \end{column}
\end{columns}
\vspace{1mm}
{\tiny\textcolor{Muted}{引用依据：引用依据为论文原文、官方规范或公开材料；本页结论按 E1 / 同行评审改进证据 的证据等级使用，不外推到未验证平台。}}
\end{frame}


\begin{frame}[t,shrink=6]{RISC-V TEE lineage: Sanctum / Keystone / Penglai / SPEAR-V：方向总结}
\DeckKicker{方向总结}
\SlideMeta{证据等级}{方向综述}{来源状态：公开材料综合}
{\large\bfseries\color{Ink} RISC-V TEE lineage: Sanctum / Keystone / Penglai / SPEAR-V 的结论是：三篇主讲材料共同给出技术演进，但仍要用证据等级约束每个结论。\par}\vspace{1.8mm}
\begin{columns}[T,totalwidth=\textwidth]
  \begin{column}{0.45\textwidth}
    \fcolorbox{Rule}{Paper}{%
  \begin{minipage}[t]{0.97\linewidth}
    \vspace{1.1mm}
    \hspace{1.4mm}\begin{minipage}{0.92\linewidth}
      \PanelTitle{内容说明}
      \scriptsize \begin{itemize}
  \item Keystone: 开源 RISC-V TEE 框架：开创/基础入口，E1 / 基础证据。
  \item Penglai: scalable enclave memory protection：代表性改进，E1 / 同行评审改进证据。
  \item SPEAR-V: low-overhead RISC-V enclave primitive：当前 SOTA/补充边界，E1 / 同行评审改进证据。
  \item 本方向结论：先用三篇主讲材料建立演进关系，再用证据等级控制机制、实验和商业化结论。
\end{itemize}
    \end{minipage}
    \vspace{1mm}
  \end{minipage}}
  \end{column}
  \begin{column}{0.49\textwidth}
    \fcolorbox{Rule}{Panel}{%
  \begin{minipage}[t]{0.97\linewidth}
    \vspace{1.1mm}
    \hspace{1.4mm}\begin{minipage}{0.92\linewidth}
      \PanelTitle{本方向方法对比}
      \scriptsize {\tiny
\begin{tabularx}{\linewidth}{@{}YYY@{}}
\toprule
\textbf{论文} & \textbf{定位} & \textbf{选择理由} \\
\midrule
Keystone: 开源 RISC-V TEE 框架 & 开创/基础入口 & EuroSys 2020 Keystone is the foundational public RISC-V open-source TEE f。 \\
Penglai: scalable enclave memory protection & 代表性改进 & OSDI 2021 Penglai is a peer-reviewed SOTA scalable RISC-V enclave OS afte。 \\
SPEAR-V: low-overhead RISC-V enclave primitive & 当前 SOTA/补充边界 & SPEAR-V is a peer-reviewed SOTA RISC-V TEE system with SoC and remote-att。 \\
\bottomrule
\end{tabularx}
}
    \end{minipage}
    \vspace{1mm}
  \end{minipage}}
  \end{column}
\end{columns}
\vspace{1mm}
{\tiny\textcolor{Muted}{引用依据：总结页综合三篇主讲材料的选择理由、评价字段和证据等级。}}
\end{frame}
